\documentclass{beamer}
\usetheme{Madrid}
\usepackage{amsmath, amssymb, listings, xcolor}

\title{Understanding Pipelining Hazards}
\subtitle{BLG 483E -- Artificial Intelligence}
\author{salihsefer36}
\date{\today}

\begin{document}

% --- Slide 1: Title ---
\begin{frame}
  \titlepage
\end{frame}

% --- Slide 2: Introduction to CPU Pipelining ---
\begin{frame}{Introduction to CPU Pipelining}
  \begin{itemize}
    \item Instruction-level parallelism (ILP) is exploited via pipelining.
    \item A pipeline divides instruction execution into $k$ discrete stages:
          IF, ID, EX, MEM, WB.
    \item Throughput approaches 1 instruction per clock cycle (IPC $\approx$ 1)
          under ideal conditions.
    \item Hazards are structural, temporal, or data-flow conflicts that prevent
          the pipeline from advancing at every clock edge.
    \item Hazard penalties manifest as pipeline stalls (bubbles) or incorrect
          architectural state.
  \end{itemize}
\end{frame}

% --- Slide 3: Data Hazards ---
\begin{frame}{Data Hazards}
  \begin{itemize}
    \item Data hazards arise from read-after-write (RAW), write-after-read (WAR),
          and write-after-write (WAW) dependences between instructions.
    \item RAW (true dependence): a consumer instruction attempts to read a
          register operand before the producer instruction completes its WB stage.
    \item WAR (anti-dependence): a later instruction writes a destination before
          a prior instruction finishes reading the same register.
    \item WAW (output dependence): two instructions write the same architectural
          register; the second write must dominate the first in program order.
    \item Detection requires a hazard detection unit that monitors the ID/EX,
          EX/MEM, and MEM/WB pipeline registers on every cycle.
    \item Mitigation: stall insertion (pipeline interlock) or operand forwarding
          (bypassing) from EX/MEM or MEM/WB to the EX stage ALU inputs.
  \end{itemize}
\end{frame}

% --- Slide 4: Structural Hazards ---
\begin{frame}{Structural Hazards}
  \begin{itemize}
    \item Structural hazards occur when multiple instructions simultaneously
          require the same hardware resource.
    \item Classic example: a unified memory with no instruction/data cache
          separation forces a stall when MEM and IF stages both require access.
    \item Register-file write-port conflicts arise when two instructions attempt
          to write-back concurrently in single-port implementations.
    \item In floating-point pipelines, multi-cycle functional units (dividers,
          square-root units) can create write-back contention.
    \item Resolution strategies: resource duplication, resource time-sharing,
          or pipeline stage reordering with compiler scheduling.
  \end{itemize}
\end{frame}

% --- Slide 5: Control Hazards ---
\begin{frame}{Control Hazards}
  \begin{itemize}
    \item Control hazards (branch hazards) arise because branch outcome and
          target address are resolved late in the pipeline (EX or MEM stage).
    \item In a 5-stage pipeline, a taken branch incurs a 2-cycle penalty if
          no mitigation is applied (instructions at IF and ID must be flushed).
    \item Static prediction strategies: predict-not-taken, predict-taken,
          BTFN (backward-taken, forward-not-taken).
    \item Dynamic prediction: 1-bit and 2-bit saturating counters in a Branch
          Prediction Buffer (BPB) indexed by the lower bits of the PC.
    \item Branch Target Buffer (BTB) stores predicted target addresses to
          eliminate target-address computation latency.
    \item Delayed branching exposes one branch delay slot filled by the compiler
          with a useful instruction or a \texttt{NOP}.
  \end{itemize}
\end{frame}

% --- Slide 6: Hazard Mitigation Summary ---
\begin{frame}{Hazard Mitigation Summary}
  \begin{itemize}
    \item \textbf{Stall (interlock):} The hazard detection unit asserts a stall
          signal, holding PC and IF/ID registers constant while inserting a bubble
          (NOP) into the ID/EX register.
    \item \textbf{Forwarding (bypassing):} Multiplexers route the result from the
          EX/MEM or MEM/WB register directly to the ALU input, reducing RAW stalls
          from 2 cycles to 0 for most cases.
    \item \textbf{Load-use hazard:} A load followed immediately by a dependent
          instruction cannot be fully resolved by forwarding alone; one stall cycle
          remains unavoidable without out-of-order execution.
    \item \textbf{Branch delay slots:} RISC ISAs (MIPS, SPARC) expose one slot
          post-branch; the instruction in this slot always executes.
    \item \textbf{Compiler scheduling:} Static reordering eliminates hazards at
          compile time by inserting independent instructions between dependent pairs.
  \end{itemize}
\end{frame}

% --- Slide 7: References ---
\begin{frame}{References}
  \begin{thebibliography}{9}
    \bibitem{patterson} Patterson, D.\,A. and Hennessy, J.\,L.
      \textit{Computer Organization and Design: The Hardware/Software Interface},
      5th ed. Morgan Kaufmann, 2014.
    \bibitem{hennessy} Hennessy, J.\,L. and Patterson, D.\,A.
      \textit{Computer Architecture: A Quantitative Approach},
      6th ed. Morgan Kaufmann, 2017.
  \end{thebibliography}
\end{frame}

\end{document}
